\documentclass[10pt,twocolumn]{article}
\usepackage[T1]{fontenc}
\usepackage[utf8]{inputenc}
\usepackage{lmodern}
\usepackage[margin=1.8cm,top=2.4cm,bottom=2cm,columnsep=0.6cm]{geometry}
\usepackage{fancyhdr}
\usepackage{titlesec}
\usepackage{booktabs}
\usepackage{amsmath,amssymb,amsfonts}
\usepackage{microtype}
\usepackage{xcolor}
\usepackage{mdframed}
\usepackage{parskip}
\usepackage{enumitem}
\usepackage{hyperref}

\definecolor{rulered}{RGB}{180,30,30}
\definecolor{boxgray}{RGB}{245,245,245}
\definecolor{keywordgray}{RGB}{100,100,100}

\pagestyle{fancy}
\fancyhf{}
\fancyhead[L]{\textsc{\small Claude Mag}}
\fancyhead[C]{\small\textit{Machine Learning Research Digest}}
\fancyhead[R]{\small 2026-09-28}
\fancyfoot[C]{\small\thepage}
\renewcommand{\headrulewidth}{0.8pt}

\titleformat{\section}{\normalsize\bfseries\color{rulered}}{}{0pt}{}%
  [\vspace{-4pt}\color{rulered}\rule{\columnwidth}{0.4pt}]
\titlespacing{\section}{0pt}{8pt}{4pt}

\hypersetup{colorlinks=true,urlcolor=rulered,linkcolor=black}

\begin{document}
\setlength{\parindent}{0pt}
\setlength{\parskip}{4pt}

{\small\color{keywordgray}\textbf{Keywords:}
  mechanistic interpretability \textbullet{}
  hallucination \textbullet{}
  sparse probing \textbullet{}
  neuron localization}\par\vspace{4pt}

{\LARGE\bfseries\raggedright
  Hallucination Neurons and Where to Find Them}\par\vspace{4pt}

{\large\itshape\raggedright
  An Investigation into the Existence of Hallucination Neurons in LLMs
  via Diagnostic Protocol}\par\vspace{6pt}

{\small\textbf{By} (authors not specified in pre-print metadata)}\par\vspace{2pt}

{\small\color{keywordgray}arXiv:2609.29781 \textbullet{} September 2026}
\par\vspace{2pt}\noindent{\color{rulered}\rule{\columnwidth}{0.6pt}}\vspace{6pt}

A five-step diagnostic protocol reveals that 19 of 22 claimed hallucination
neurons in Gemma~3~4B have Pearson $|r|>0.7$ with other features, exhibit
only moderate bootstrap stability, and show weak sparse-dense ranking overlap
--- demonstrating that sparse predictive structure coexists with non-unique
neuron selection, and that routine diagnostic validation is necessary before
making causal localization claims.

\begin{mdframed}[backgroundcolor=boxgray,linecolor=rulered,linewidth=1pt,
    innerleftmargin=6pt,innerrightmargin=6pt,
    innertopmargin=6pt,innerbottommargin=6pt]
\textbf{\small AT A GLANCE}\par\vspace{4pt}
\begin{description}[leftmargin=0pt,labelsep=4pt,itemsep=2pt,parsep=0pt,
    font=\small\bfseries]
  \item[Problem:] Sparse probing claims specific neurons causally drive LLM
    hallucination, but these claims are not tested against known failure modes
    of L1-regularized probing in correlated feature spaces.
  \item[Why it matters:] Non-unique neuron selection undermines the validity
    of targeted steering and auditing based on claimed hallucination neurons.
  \item[Method:] A five-step diagnostic protocol: feature correlation,
    bootstrap stability, sparse vs.\ dense ranking, intervention baselines,
    and cross-dataset evaluation.
  \item[Result:] 19/22 H-Neurons fail the feature correlation criterion in
    Gemma~3~4B; bootstrap stability and sparse-dense agreement are weak.
  \item[Evidence:] TriviaQA, BioASQ, and NQ-Open across Gemma~3~4B and
    MedGemma~4B.
\end{description}
\end{mdframed}

\section{The Big Picture}

Mechanistic interpretability of large language models has made rapid progress
in identifying circuits, features, and neurons associated with specific
behaviors. A particularly consequential line of work claims to identify
``hallucination neurons''---small sets of neurons whose activation levels
predict and causally drive factual hallucinations. If such neurons exist and
are uniquely identifiable, they offer a compelling target for behavioral
steering: suppress the neuron, suppress the hallucination.

But neuron identification typically relies on L1-regularized (Lasso-style)
sparse probing, which is well-known to select non-unique solutions in the
presence of correlated features. A Lasso probe that achieves high AUROC does
not establish that the selected neurons are the unique causal locus of the
behavior---only that they are predictive given the regularization and the
training data. In high-dimensional, highly correlated neural activation
spaces, this distinction matters enormously.

This paper proposes a minimum diagnostic standard for neuron-localization
claims and applies it to the H-Neuron literature, revealing systematic
non-uniqueness in the claimed hallucination neurons.

\section{Why Existing Approaches Fall Short}

\textbf{L1-regularized probing} selects solutions non-uniquely when features
are correlated: many alternative neuron sets achieve similar or better
discrimination. The selected set depends on regularization strength,
initialization, and the specific training instances.

\textbf{High AUROC} is a necessary but not sufficient condition for
localization. Many overlapping neuron subsets may achieve high AUROC; the
probe's predictive success does not establish which neurons are causally
necessary.

\textbf{Ablation studies} (removing a selected neuron and measuring behavior
change) are often compared against no-ablation baselines rather than against
equally predictive alternative neurons, making causal claims fragile.

\textbf{Single-dataset evaluation} cannot distinguish neurons that detect
hallucination generically from neurons whose predictive power is specific to
the evaluation dataset's format or domain.

\section{Core Mechanism}

\textbf{Five-Step Diagnostic Protocol.} The paper proposes five diagnostic
criteria that should be satisfied before claiming unique neuron localization:

\begin{enumerate}[itemsep=2pt,parsep=0pt]
  \item \textbf{Feature correlation:} Flag selected neurons with Pearson
    $|r|>0.7$ with unselected features as potentially non-unique.
  \item \textbf{Bootstrap stability:} Measure Jaccard overlap of the selected
    neuron set across bootstrap samples of the probe training data.
  \item \textbf{Sparse vs.\ dense ranking:} Compare the ranking of neurons
    by the L1 probe with their ranking by a dense (non-regularized) importance
    score; low overlap signals L1-driven selection, not importance.
  \item \textbf{Intervention baselines:} Compare ablation of selected neurons
    against ablation of equally predictive randomly drawn neurons; a uniquely
    causal neuron should show strictly greater behavior change.
  \item \textbf{Cross-dataset evaluation:} Evaluate selected neurons on
    held-out datasets not used for selection; poor transfer indicates
    dataset-specific rather than general hallucination neurons.
\end{enumerate}

\section{Technical Formulation}

Let $\mathbf{h} \in \mathbb{R}^d$ be activations at a layer, $y \in \{0,1\}$
the hallucination label. The L1-regularized probe solves:
\[
\hat{\beta} = \arg\min_\beta \;\mathcal{L}(y,\, \mathbf{h}^\top\beta)
  + \lambda \|\beta\|_1
\]
Selected neuron set: $S = \{i : \hat{\beta}_i \neq 0\}$.

Criterion 1 (feature correlation): Flag $S$ if
$\exists\, j \notin S, i \in S$ such that $|\text{Corr}(h_i, h_j)| > 0.7$.

Criterion 2 (bootstrap stability): $J_B = |S \cap S_b| / |S \cup S_b|$
for bootstrap sample $b$; flag if $\mathbb{E}[J_B] < 0.8$.

Criterion 3 (ranking disagreement): Spearman $\rho(\hat{\beta},
\tilde{\beta})$ where $\tilde{\beta}$ is the dense-probe importance;
flag if $\rho < 0.5$.

Cross-dataset AUROC gaps (Gemma~3~4B vs MedGemma~4B):
\[
\Delta_\text{TriviaQA} = +0.311 \;\text{vs}\; +0.235
\]
\[
\Delta_\text{BioASQ} = +0.474 \;\text{vs}\; +0.455
\]
\[
\Delta_\text{NQ-Open} = +0.128 \;\text{vs}\; +0.112
\]

\section{Learning or Inference Procedure}

The diagnostic is applied post-hoc to trained probes:
\begin{enumerate}[itemsep=2pt,parsep=0pt]
  \item Train an L1 probe on each (model, dataset) pair.
  \item Record the selected neuron set $S$.
  \item Apply each diagnostic criterion.
  \item Report pass/fail and overall diagnostic verdict.
\end{enumerate}
No model training is performed; the protocol is a validation framework.

\section{What the Guarantee Says}

When all five criteria are satisfied, the selected neuron set is consistent
with being a unique causal locus of the behavior. When criteria 1--3 fail
(as they do here), the probe's predictive power is consistent with a broader
class of alternative neuron sets, and the causal localization claim is not
established. This does not falsify hallucination detection; it shows the
detection is distributed, not localized.

\section{Experimental Findings}

\begin{itemize}[itemsep=2pt,parsep=0pt]
  \item \textbf{19 of 22} H-Neurons in Gemma~3~4B have Pearson $|r|>0.7$
    with at least one unselected feature---failing criterion~1
  \item Bootstrap Jaccard overlap is moderate, confirming sensitivity to
    training subsample---failing criterion~2
  \item Sparse-dense rank correlation is weak, indicating L1 regularization
    drives the selection---failing criterion~3
  \item Gemma~3~4B consistently outperforms MedGemma~4B on hallucination
    probe AUROC across all three datasets
  \item Cross-dataset probe quality varies, suggesting partial dataset
    specificity in neuron selection
\end{itemize}

\section{Ablations and Interpretation}

\textbf{AUROC vs.\ uniqueness:} High AUROC is achieved by many overlapping
neuron sets; the reported probe AUROC does not support localization.

\textbf{Dataset specificity:} Different datasets select partially different
neuron sets, consistent with distributed rather than universal hallucination
neurons.

\textbf{Implications for steering:} Interventions on L1-selected neurons may
or may not generalize because many correlated alternatives exist. Robust
steering likely requires denser probes, ensemble methods, or causal
identification beyond predictive accuracy.

\textbf{Broader message:} Routine diagnostic validation---particularly feature
correlation and bootstrap stability checks---should be a minimum requirement
for sparse-neuron localization claims in interpretability research.

\section*{Reference}

\textbf{Hallucination Neurons and Where to Find Them: An Investigation
into the Existence of Hallucination Neurons}.
arXiv:2609.29781, September 2026.\\
\url{https://arxiv.org/abs/2609.29781}

\end{document}
